\documentclass[a4paper,11pt]{scrartcl}
% preamble.tex -- ALL typography for the manuscript lives here.
%
% Changelog (reverse-chronological)
%   2026-09-01  arXiv build.  (a) Dropped "number within=section" from the
%               theorem counter: Typst numbers ALL theorem-like environments
%               globally and continuously (Theorem 1, Remark 2, Definition 3,
%               ...), and the section-scoped LaTeX numbering silently produced a
%               different sequence from the one the article's own prose, the
%               Lean sources and the project notes all cite.
%               (b) Added tikz (for the one hand-drawn figure, which pandoc
%               cannot translate) and a heading-less thebibliography, because
%               the body already carries its own numbered "References" section
%               heading converted from the Typst.
%   2026-08-30  Created, porting the Typst theorion/cosmos.fancy look to LaTeX.
%
% The point of this file, per Jan: submitting to a journal should be a ONE-FILE
% swap.  Nothing about fonts, colours, box geometry or theorem styling appears
% in the manuscript body -- the body uses only \begin{theorem}...\end{theorem}
% and friends, which every journal class already provides.
%
% To retarget:  set \documentclass to the journal's, delete the tcolorbox block
% below, and the body compiles unchanged with the class's own theorem styles.

% ---- class -------------------------------------------------------------
% KOMA (scrartcl) rather than amsart: better control of headings and margins,
% and it is not a journal class, so nothing here pretends to be a submission
% format.  An amsart variant is kept at preamble-amsart.tex.
\usepackage[a4paper,margin=2.6cm]{geometry}
\usepackage{amsmath,amssymb,amsthm,mathtools}
\usepackage{fontspec}                 % xelatex
\usepackage{microtype}
\usepackage[hidelinks]{hyperref}
\usepackage[capitalise,noabbrev]{cleveref}
\usepackage{booktabs}
% pandoc emits longtable for every table, plus these helpers; without them the
% build dies with "Environment longtable undefined" and a cascade of
% "Misplaced \noalign".
\usepackage{longtable,array,calc}
\usepackage{graphicx}
% Converted prose is dense with unbreakable inline maths, which leaves TeX no
% legal break in some lines; pandoc's own default template sets this for the
% same reason.  Without it, a handful of lines run ~25pt into the margin.
\setlength{\emergencystretch}{3em}
\providecommand{\tightlist}{\setlength{\itemsep}{0pt}\setlength{\parskip}{0pt}}
\providecommand{\pandocbounded}[1]{#1}
\usepackage[most]{tcolorbox}
\tcbuselibrary{theorems}
\usepackage{varwidth}                 % for tcolorbox's "varwidth boxed title"
\usepackage{tikz}
\usetikzlibrary{arrows.meta,positioning,shapes.geometric}

% ---- bibliography without its own heading ------------------------------
% The converted body already carries "\section{References}" (Typst's "=
% References", numbered 10.).  Stock thebibliography would add a SECOND,
% unnumbered "References" heading right under it.  This is the standard
% article.cls environment with the \section*{\refname} line removed.
\makeatletter
\renewenvironment{thebibliography}[1]
  {\list{\@biblabel{\@arabic\c@enumiv}}%
        {\settowidth\labelwidth{\@biblabel{#1}}%
         \leftmargin\labelwidth
         \advance\leftmargin\labelsep
         \@openbib@code
         \usecounter{enumiv}%
         \let\p@enumiv\@empty
         \renewcommand\theenumiv{\@arabic\c@enumiv}}%
   \sloppy
   \clubpenalty4000
   \@clubpenalty \clubpenalty
   \widowpenalty4000%
   \sfcode`\.\@m}
  {\def\@noitemerr
    {\@latex@warning{Empty `thebibliography' environment}}%
   \endlist}
\providecommand{\@openbib@code}{}
\makeatother

% ---- the coloured boxed theorem style ---------------------------------
% NOTE the empty {} prefix on every \newtcbtheorem below.  It does NOT make the
% labels bare -- tcolorbox prepends a colon unconditionally, so
% \begin{definition}{t}{def-y1y2} defines ":def-y1y2" (checked in main.aux, not
% assumed).  An empty prefix only shortens "prp:foo" to ":foo".  Cross-references
% are therefore written by tools/markers2env.py as
% \hyperref[:foo]{Definition~\ref*{:foo}}, which also supplies the environment
% name -- every box here shares one counter, so cleveref could only ever call
% all of them "Theorem".  See fix_references' docstring for the full argument.
% Reproduces theorion's cosmos.fancy: a tinted body, a saturated title bar,
% rounded corners, a coloured left rule.
% An attached boxed title is a single unbreakable line by default, so a long
% Typst title just runs off the paper: Remark 27's ("... component-by-component
% is what Proposition 26 answers") was cut at "is what P", silently -- the rest
% is in the page stream, outside the trim.  varwidth lets it wrap.
\tcbset{varwidth boxed title*=-1cm}

\definecolor{clrThm}{HTML}{2F6FAF}
\definecolor{clrDef}{HTML}{1F7A5A}
\definecolor{clrRem}{HTML}{8A6D1F}
\definecolor{clrExa}{HTML}{6B4C9A}

% NO "number within=section": Typst numbers every theorem-like environment from
% one global, continuous counter (Theorem 1, Remark 2, Definition 3, Example 4,
% Problem 5, Theorem 6, ...), and those numbers are cited by name elsewhere.
\newtcbtheorem{theorem}{Theorem}{%
  enhanced, breakable, colback=clrThm!4, colframe=clrThm, coltitle=white,
  fonttitle=\bfseries, arc=2pt, boxrule=0.6pt, left=6pt, right=6pt,
  attach boxed title to top left={xshift=6pt,yshift=-2mm},
  boxed title style={colback=clrThm, arc=1pt}}{}

% proposition/lemma share the theorem colour and counter
\newtcbtheorem[use counter from=theorem]{proposition}{Proposition}{%
  enhanced, breakable, colback=clrThm!4, colframe=clrThm, coltitle=white,
  fonttitle=\bfseries, arc=2pt, boxrule=0.6pt, left=6pt, right=6pt,
  attach boxed title to top left={xshift=6pt,yshift=-2mm},
  boxed title style={colback=clrThm, arc=1pt}}{}
\newtcbtheorem[use counter from=theorem]{lemma}{Lemma}{%
  enhanced, breakable, colback=clrThm!4, colframe=clrThm, coltitle=white,
  fonttitle=\bfseries, arc=2pt, boxrule=0.6pt, left=6pt, right=6pt,
  attach boxed title to top left={xshift=6pt,yshift=-2mm},
  boxed title style={colback=clrThm, arc=1pt}}{}
\newtcbtheorem[use counter from=theorem]{definition}{Definition}{%
  enhanced, breakable, colback=clrDef!4, colframe=clrDef, coltitle=white,
  fonttitle=\bfseries, arc=2pt, boxrule=0.6pt, left=6pt, right=6pt,
  attach boxed title to top left={xshift=6pt,yshift=-2mm},
  boxed title style={colback=clrDef, arc=1pt}}{}
\newtcbtheorem[use counter from=theorem]{remark}{Remark}{%
  enhanced, breakable, colback=clrRem!4, colframe=clrRem, coltitle=white,
  fonttitle=\bfseries, arc=2pt, boxrule=0.6pt, left=6pt, right=6pt,
  attach boxed title to top left={xshift=6pt,yshift=-2mm},
  boxed title style={colback=clrRem, arc=1pt}}{}
\newtcbtheorem[use counter from=theorem]{example}{Example}{%
  enhanced, breakable, colback=clrExa!4, colframe=clrExa, coltitle=white,
  fonttitle=\bfseries, arc=2pt, boxrule=0.6pt, left=6pt, right=6pt,
  attach boxed title to top left={xshift=6pt,yshift=-2mm},
  boxed title style={colback=clrExa, arc=1pt}}{}

\definecolor{clrPrb}{HTML}{A0522D}
\newtcbtheorem[use counter from=theorem]{problem}{Problem}{%
  enhanced, breakable, colback=clrPrb!4, colframe=clrPrb, coltitle=white,
  fonttitle=\bfseries, arc=2pt, boxrule=0.6pt, left=6pt, right=6pt,
  attach boxed title to top left={xshift=6pt,yshift=-2mm},
  boxed title style={colback=clrPrb, arc=1pt}}{}

% proof stays plain -- a boxed proof reads badly at this density
\renewcommand{\qedsymbol}{$\square$}

% ---- cross-references into the boxes -----------------------------------
% \thmref{Definition}{def-y1y2}  ->  "Definition 14", hyperlinked.
% The supplement is passed in because every box above shares ONE counter
% (tcb@cnt@theorem), so cleveref has no way to tell a Lemma from a Remark; and
% the ":" is tcolorbox's own, prepended to the third argument regardless of the
% \newtcbtheorem prefix.
% Deliberately NO square brackets in the call: an earlier version emitted
% \hyperref[:foo]{...} directly, and amsthm builds its proof heading as
% \item[...#1...], whose optional-argument scanner ends at the first "]" -- the
% one inside \hyperref's own [..].  204 errors, all from that.
\newcommand{\thmref}[2]{\hyperref[:#2]{#1~\ref*{:#2}}}

% ---- displays that are wider than the measure --------------------------
% Typst shrinks an over-wide display to fit (Example 38's 3x9 Jacobian, split
% into a 3x6 block, is typeset visibly smaller there).  LaTeX instead lets it
% run off the paper: 144pt of the matrix was outside the trim, with only an
% Overfull \hbox in the log to say so.  \fitdisplay measures first and scales
% ONLY when the content really is too wide, so ordinary displays are untouched.
\makeatletter
\newsavebox{\jts@fitbox}
\newcommand{\fitdisplay}[1]{%
  \sbox{\jts@fitbox}{\ensuremath{\displaystyle #1}}%
  \ifdim\wd\jts@fitbox>\linewidth
    \resizebox{\linewidth}{!}{\usebox{\jts@fitbox}}%
  \else
    \usebox{\jts@fitbox}%
  \fi}
\makeatother

% ---- notation ----------------------------------------------------------
\newcommand{\FF}{\mathbb{F}}
\newcommand{\QQ}{\mathbb{Q}}
\newcommand{\RR}{\mathbb{R}}
\newcommand{\AAA}{\mathbb{A}}
\newcommand{\PP}{\mathbb{P}}
\newcommand{\vek}[1]{\overline{#1}}

% Filenames like code/verify\_generator\_independence.sage occur in running
% prose and have no legal break point, so TeX pushes 148pt of one into the
% margin.  Allow a break after each underscore; \allowbreak sets nothing, so
% the filename is unchanged, it just becomes breakable.
\let\jtsOldUnderscore\_
\renewcommand{\_}{\jtsOldUnderscore\allowbreak}
\newcommand{\enc}{\operatorname{enc}}
\newcommand{\red}{\operatorname{red}}
\DeclareMathOperator{\tr}{tr}
\DeclareMathOperator{\pd}{pd}

\begin{document}
\section{Introduction}

\(2 \times 2\) matrices \(A,B\) have commutator \(\lbrack A,B\rbrack = AB - BA\), and the \textbf{iterated} left-normed commutator of \(k\) matrices is \(\left\lbrack A_{1},\ldots,A_{k} \right\rbrack ≔ \left\lbrack \ldots\left\lbrack \left\lbrack A_{1},A_{2} \right\rbrack,A_{3} \right\rbrack,\ldots,A_{k} \right\rbrack\), matching the standard lower-central-series convention for Lie algebras. This object arose out of a question in algebraic complexity theory -- how many scalar multiplications are needed to \textbf{compute} \(\left\lbrack A_{1},\ldots,A_{k} \right\rbrack\)? -- which remains genuinely open for \(k \geq 3\) and is treated in full in the companion note k-iterated-commutators-complexity.typ. The one fact from that story needed here is purely algebraic, not complexity-theoretic: an explicit bilinear formula for the ordinary \(2 \times 2\) commutator, chained across \(k\) matrices, gives an explicit polynomial parametrization of \(\left\lbrack A_{1},\ldots,A_{k} \right\rbrack\)`s reduced coordinates (\cref{:def-wk} below), and it is the \textbf{vanishing locus} of this parametrization, as an algebraic variety and its coordinate ring, that is the actual subject of this article. (One naming note: the filename this article is kept under, ring-of-k-commuting-2x2-matrices, uses ``commuting'' in this generalized iterated-commutator-vanishing sense -- \(\left\lbrack A_{1},\ldots,A_{k} \right\rbrack = 0\) -- not literally ``every pair of \(A_{1},\ldots,A_{k}\) commutes,'' a weaker condition than the classical pairwise commuting variety it specializes to at \(k = 2\).)

\textbf{The classical commuting variety.} For \(A,B\) arbitrary (not necessarily \(2 \times 2\)) \(n \times n\) matrices, the variety \(C_{n} ≔ \left\{ (A,B) \in M_{n} \times M_{n}:AB = BA \right\}\) -- the ordinary, unreduced, un-iterated case this article's own \(k = 2\) specializes to -- has been studied since the 1950s. Its irreducibility was proven independently by Motzkin and Taussky MotzkinTaussky1955 and by Gerstenhaber Gerstenhaber1961, the latter also establishing the now-standard bound \(\dim_{\mathbb{F}}\mathbb{F}\lbrack A,B\rbrack \leq n\) for the subalgebra generated by two commuting matrices. Whether the coordinate ring \(\mathbb{F}\lbrack A,B\rbrack/I\), \(I\) the ideal generated by the entries of \(AB - BA\), is \textbf{reduced} and \textbf{Cohen-Macaulay} for every \(n\) is a much harder question, attributed to Michael Artin and Melvin Hochster -- an unpublished, informally-circulated conjecture with no single dateable original source ArtinHochsterConjecture, stated precisely as ``Conjecture 1.1'' in Majidi-Zolbanin and Snapp MajidiZolbaninSnapp. (A closely related, more specifically dated 1984 reducedness conjecture, for a different but adjacent ideal built from a commutator matrix's diagonal/anti-diagonal entries, is attributed to Hochster alone by Chau Chau2024.) Freyja Hreinsdóttir's body of work Hreinsdottir2005 -- building on her own 1997 Stockholm doctoral thesis -- attacks the Artin-Hochster conjecture case by case at \(n = 3,4\): syzygies, Betti numbers, the canonical module, and a Koszul-dual structure theorem, verified computationally through \(n = 7\) but proven in general only in these small cases. Knutson Knutson2003 studies the minimal primes of closely related schemes from a different angle. Cohen-Macaulayness of the \textbf{classical} commuting variety therefore remains open, for every \(n \geq 3\), at the time of writing.

At \(n = 2\) specifically -- this article's \(k = 2\) base case -- the question is comparatively easy, and is in effect answered again, as a special case, by \cref{:thm-cm} below: \(I_{2}\) (\cref{:def-wk}) is literally the ideal of the \(n = 2\) commuting variety, expressed in reduced (trace-removed) coordinates rather than the full matrix-entry coordinates the classical literature typically uses. The genuinely new object in this article is the \textbf{iterated} generalization to \(k \geq 3\): not the classical commuting variety at all (see the naming note above), but a structurally analogous determinantal ideal, built by chaining the same bilinear cross-product formula that governs the \(n = 2\) case. It is exactly the same \textbf{flavor} of question the Artin-Hochster conjecture asks -- is this natural, geometrically meaningful ideal Cohen-Macaulay? is it reduced? -- for a different, recursively-defined family of ideals, and this article answers it completely and unconditionally, for every \(k \geq 2\) (\cref{:thm-primarydecomp}), rather than case by case or conjecturally. In that sense the results below can be read as a fully worked structural analogue of the still-open classical question, in a setting simple enough (\(n = 2\), but \(k\) unbounded) to admit a complete answer.

\section{Notation}

Notation used across more than one later result, in order of first appearance; symbols local to a single proof are omitted.

\textbf{A typographic convention, used throughout}: an \textbf{overline} denotes a \textbf{vector} -- an element of \(\mathbb{F}^{3}\) (or, for \(R_{k}\) itself, of the free module \(R_{k}^{3}\)) -- while the \textbf{same letter, plain, with a two-index comma subscript}, denotes one of that vector's three individual scalar \textbf{components}. So \({\overline{v}}_{i} \in \mathbb{F}^{3}\) is block \(i\)`s whole reduced-coordinate vector, while its three components are the plain scalars \(v_{i,1},v_{i,2},v_{i,3} \in \mathbb{F}\) (equivalently \(v_{i,j}\) for \(j = 1,2,3\)) -- never overlined, since a single scalar entry is not itself a vector. The same convention applies to \({\overline{w}}_{j}\) versus its components \(w_{j,1},w_{j,2},w_{j,3}\).

\begin{longtable}[]{@{}cll@{}}
\toprule\noalign{}
\endhead
\bottomrule\noalign{}
\endlastfoot
\textbf{Symbol} & \textbf{Meaning} & \textbf{Defined in} \\
\(\overline{v}\) (also \({\overline{v}}_{i}\); components \(v_{i,1},v_{i,2},v_{i,3}\), or \(v_{i,j}\)) & reduced-coordinate vector in \(\mathbb{F}^{3}\), and its blocks & §3, Setup \\
\(\tau\) (and its matrix \(T\)) & the reshuffle \(\tau(a,b,c) ≔ (b,a,2c)\) & §3, Setup (\(T\): \cref{:lem-jacrecursion}) \\
\(\mathbb{F}\) & ambient field, characteristic \(\neq 2\) & §3, after \cref{:def-wk}`s example \\
\(R_{k}\) & the polynomial ring \(\mathbb{F}\left\lbrack {\overline{v}}_{1},\ldots,{\overline{v}}_{k} \right\rbrack\) & \cref{:def-wk} \\
\(I_{k}\) & the defining ideal, \(I_{k} \subset R_{k}\) & \cref{:def-wk} \\
\({\overline{w}}_{j}\) & (a) the recursive vector in \(\mathbb{F}^{3}\); (b) as a polynomial map \(\mathbb{A}^{3j} \rightarrow \mathbb{A}^{3}\) & (a) \cref{:def-wk}; (b) \cref{:lem-nondeg} \\
\(U_{j}\) & the Zariski-open locus where \({\overline{w}}_{j} \neq 0\) & \cref{:lem-nondeg} \\
\(\pi\) & the projection \(\mathbb{A}^{3k} \rightarrow \mathbb{A}^{3(k - 1)}\) forgetting \({\overline{v}}_{k}\) & \cref{:def-y1y2} \\
\(\Phi\) & the map \(\mathbb{A}^{3(k - 1)} \times \mathbb{A}^{1} \rightarrow \mathbb{A}^{3k}\) & \cref{:def-y1y2} \\
\(Y_{1}\), \(Y_{2}\) & the two strata of \(V\left( I_{k} \right)\) & \cref{:def-y1y2} \\
\(P_{j}\) (\(j = 2,\ldots,k\)) & the primary components of \(I_{k}\) & \cref{:def-primecomponents} \\
\(J_{k}\) & the Jacobian of \({\overline{w}}_{k}\) & \cref{:lem-jacrecursion} \\
\(\lbrack u\rbrack_{\times}\) & skew matrix with \(\lbrack u\rbrack_{\times}v = u \times v\) & \cref{:lem-jacrecursion} \\
\end{longtable}

\section{Setup: reduced coordinates and the defining ideal}

\textbf{Step 1: discard the scalar part.} Write \(A = \begin{pmatrix}
a_{11} & a_{12} \\
a_{21} & a_{22}
\end{pmatrix}\) as the sum of its scalar part \(\left( \text{tr}(A)/2 \right)I\) and its traceless remainder \(A_{0} ≔ A - \left( \text{tr}(A)/2 \right)I = \begin{pmatrix}
d_{a}/2 & a_{12} \\
a_{21} & - d_{a}/2
\end{pmatrix}\), \(d_{a} ≔ a_{11} - a_{22}\). Scalar matrices are central, so they never contribute to a commutator: \(\lbrack cI,B\rbrack = cIB - BcI = cB - cB = 0\) for every \(B\) and every scalar \(c\), hence \(\lbrack A,B\rbrack = \left\lbrack A_{0},B_{0} \right\rbrack\) identically -- no information relevant to \textbf{any} bracket is lost by discarding the scalar part, so from here on only the traceless remainder matters.

\textbf{Step 2: encode the traceless remainder as a triple.} A traceless matrix \(\begin{pmatrix}
x & y \\
z & - x
\end{pmatrix}\) has exactly three free entries -- the \((2,2)\) entry is forced once \(x,y,z\) are fixed -- giving a bijection onto \(\mathbb{R}^{3}\). Write \[\text{ red}(A) ≔ \left( a_{12},a_{21},d_{a} \right) \in \mathbb{R}^{3}\] for this encoding (so \(\text{red}(A) = \text{ red}\left( A_{0} \right)\), reading \(A\)`s traceless part off directly without subtracting the scalar part by hand first) -- this is the notation \(\overline{v}\) used throughout the rest of the article. Matching Step 1's \(A_{0} = \begin{pmatrix}
d_{a}/2 & a_{12} \\
a_{21} & - d_{a}/2
\end{pmatrix}\): \(\text{red}\)`s three coordinates are literally \(A_{0}\)`s free entries, except that \(d_{a}\) stands in for \textbf{twice} the \((1,1)\)-entry rather than the entry itself -- a convention fixed once and for all here purely so that Step 4 below comes out with no stray constants.

\textbf{Step 3: form the commutator.} Given traceless \(A_{0} = \begin{pmatrix}
x & y \\
z & - x
\end{pmatrix}\) and \(B_{0} = \begin{pmatrix}
u & v \\
w & - u
\end{pmatrix}\), the ordinary matrix commutator \(\left\lbrack A_{0},B_{0} \right\rbrack = A_{0}B_{0} - B_{0}A_{0}\) is again traceless automatically (\(\text{tr}(PQ) = \text{tr}(QP)\) for any \(P,Q\), so \(\text{tr}\left( \lbrack P,Q\rbrack \right) = 0\) always), hence again encodable by \(\text{red}\).

\textbf{Step 4: the encoded bracket is the cross product -- up to one fixed reshuffle.} Carrying out Step 3's multiplication in coordinates (the \(\mathfrak{sl}_{2}({\mathbb{C}}) \cong \mathfrak{so}(3) \cong \left( \mathbb{R}^{3}, \times \right)\) isomorphism made explicit) and writing \(\left\lbrack A_{0},B_{0} \right\rbrack = \begin{pmatrix}
x' & y' \\
z' & - x'
\end{pmatrix}\), one checks directly that, with \(\text{red}(A)\) and \(\text{red}(B)\) standing in for themselves (\textbf{not} abbreviated to bare letters here, to avoid colliding with Step 3's plain scalars \(u,v,w\)), \[z' = \left( \text{red}(A) \times \text{ red}(B) \right)_{1},\quad y' = \left( \text{red}(A) \times \text{ red}(B) \right)_{2},\quad x' = \left( \text{red}(A) \times \text{ red}(B) \right)_{3.}\] Reading these three equalities off against \(\text{red}\left( \lbrack A,B\rbrack \right) = (y',z',2x')\) (Step 2's convention, applied to \(\lbrack A,B\rbrack\) itself) gives \[\text{ red}\left( \lbrack A,B\rbrack \right) = \tau(\text{red}(A) \times \text{ red}(B)),\quad\tau(a,b,c) ≔ (b,a,2c).\] The right-hand side is \textbf{not} the raw cross product \(\text{red}(A) \times \text{ red}(B)\) itself, but that cross product composed with the fixed invertible reshuffle \(\tau\) (invertible whenever \(2 \neq 0\)): the \((1,2)\)-swap comes from \(\text{red}\)`s coordinate order (\(a_{12}\) before \(a_{21}\), the transpose of the naive \((x,y,z)\) order), and the factor of 2 on the third slot comes from Step 2's \(d_{a} = 2x\) convention. Both are bookkeeping artifacts of the coordinate choice fixed in Step 2 -- \(\tau\) is already forced at this single-bracket level, not a new phenomenon introduced only once brackets start chaining below.

\begin{remark}{why \cref{:def-wk}`s \({\overline{w}}_{j}\) carries no outer \(\tau\)}{rem-red-tau-omitted} \cref{:def-wk} below sets \({\overline{w}}_{2} ≔ {\overline{v}}_{1} \times {\overline{v}}_{2}\) -- the \textbf{untwisted} cross product, not \(\text{red}\left( \lbrack A,B\rbrack \right) = \tau({\overline{v}}_{1} \times {\overline{v}}_{2})\) itself as Step 4 above would suggest. This is deliberate: for any fixed invertible linear map \(\tau\) and any polynomial vector \(f = \left( f_{1},f_{2},f_{3} \right)\), the ideal generated by the components of \(\tau(f)\) equals the ideal generated by the components of \(f\) -- each component of one is an \(\mathbb{F}\)-linear combination of the components of the other (via \(\tau\), respectively \(\tau^{- 1}\)), so the two generating sets differ only by an invertible change of generators, and invertible changes of generators do not change the ideal generated. Hence \(I_{2} ≔ \left( w_{2,1},w_{2,2},w_{2,3} \right)\) already equals the ideal generated by \(\text{red}\left( \lbrack A,B\rbrack \right)\)`s three components, with no need to carry the extra \(\tau\) explicitly -- \({\overline{w}}_{j}\) is used purely because it is the cheaper quantity to compute recursively (one cross product per step, not a cross product composed with \(\tau\) at every stage), and the single \(\tau\) that \textbf{does} appear in \cref{:def-wk}`s recursion below is exactly the one deferred from the previous step, not a new one introduced fresh at each stage.

\end{remark} \protect\phantomsection\label{rem-red-tau-omitted}{}

Chaining this bilinear formula across \(k\) matrices -- each new bracket stage reuses the previous stage's output coordinates directly, no new input data needed -- requires tracking how one stage's raw cross-product output is reinterpreted as the next stage's input vector: precisely the \(\tau\) just discussed, applied once per new stage. This gives the following exact recursive description:

\begin{definition}{reduced iterated commutator}{def-wk} \[{\overline{w}}_{2} ≔ {\overline{v}}_{1} \times {\overline{v}}_{2},\quad{\overline{w}}_{j} ≔ \tau({\overline{w}}_{j - 1}) \times {\overline{v}}_{j}\quad(j = 3,\ldots,k).\] Let \(R_{k} ≔ \mathbb{Q}\left\lbrack {\overline{v}}_{1},\ldots,{\overline{v}}_{k} \right\rbrack\) (\({\overline{v}}_{i} \in \mathbb{Q}^{3}\), \(3k\) variables) and \(I_{k} \subset R_{k}\) the ideal generated by the three components of \({\overline{w}}_{k}\) -- the vanishing ideal of \(\left\lbrack A_{1},\ldots,A_{k} \right\rbrack = 0\) in reduced coordinates.

\end{definition} \protect\phantomsection\label{def-wk}{}

To spell out exactly what this means: ``\(R_{k} ≔ \mathbb{Q}\left\lbrack {\overline{v}}_{1},\ldots,{\overline{v}}_{k} \right\rbrack\)'' is shorthand for \textbf{the polynomial ring in the \(3k\) scalar variables \(v_{i,j}\) (\(i = 1,\ldots,k\), \(j = 1,2,3\))}, grouped into \(k\) blocks of 3 -- \textbf{not} a ring with \(k\) vector-valued generators (a polynomial ring's generators are always scalars). At \(k = 2\), written out in full: \[R_{2} = \mathbb{Q}\left\lbrack v_{1,1},v_{1,2},v_{1,3},v_{2,1},v_{2,2},v_{2,3} \right\rbrack,\] six ordinary scalar indeterminates, with \({\overline{v}}_{1} = \left( v_{1,1},v_{1,2},v_{1,3} \right)\) and \({\overline{v}}_{2} = \left( v_{2,1},v_{2,2},v_{2,3} \right)\) merely a notational grouping of them into two blocks of three -- convenient because \({\overline{w}}_{2} = {\overline{v}}_{1} \times {\overline{v}}_{2}\) is most naturally written as a cross product of two block-vectors, not as some formula in six unrelated scalars.

Each \({\overline{w}}_{j}\) is a triple of polynomials, i.e. \({\overline{w}}_{j} \in R_{j}^{3}\) (the free rank-3 module over \(R_{j}\)) -- equivalently, exactly as recorded in the Notation table's entry (b) for \({\overline{w}}_{j}\), the polynomial map \({\overline{w}}_{j}:\mathbb{A}^{3j} \rightarrow \mathbb{A}^{3}\) \textbf{given by substitution}: evaluating each of \({\overline{w}}_{j}\)`s three polynomial components at a point \(\left( {\overline{v}}_{1},\ldots,{\overline{v}}_{j} \right) \in \mathbb{A}^{3j}\). Concretely at \(k = 2\), writing out all six variables, \({\overline{w}}_{2}\) is the substitution map \[\begin{pmatrix}
v_{1,1} \\
v_{1,2} \\
v_{1,3} \\
v_{2,1} \\
v_{2,2} \\
v_{2,3}
\end{pmatrix} \in \mathbb{A}^{6}\quad \mapsto \quad\begin{pmatrix}
v_{1,2}v_{2,3} - v_{1,3}v_{2,2} \\
v_{1,3}v_{2,1} - v_{1,1}v_{2,3} \\
v_{1,1}v_{2,2} - v_{1,2}v_{2,1}
\end{pmatrix} \in \mathbb{A}^{3},\] and \(I_{k}\) (\cref{:def-wk}) is, equivalently, the ideal generated by the three coordinate functions of \({\overline{w}}_{k}:\mathbb{A}^{3k} \rightarrow \mathbb{A}^{3}\) pulled back along this substitution.

\begin{example}{{}\({\overline{w}}_{k}\) and \(I_{k}\), explicitly, for \(k = 2,3,4\)}{thm-determinantal} Writing \({\overline{v}}_{i} = \left( v_{i,1},v_{i,2},v_{i,3} \right)\) for the three coordinates of block \({\overline{v}}_{i}\): at \(k = 2\), the three components of \({\overline{w}}_{2} = {\overline{v}}_{1} \times {\overline{v}}_{2}\) are \[v_{1,2}v_{2,3} - v_{1,3}v_{2,2},\quad v_{1,3}v_{2,1} - v_{1,1}v_{2,3},\quad v_{1,1}v_{2,2} - v_{1,2}v_{2,1},\] and \(I_{2} ≔ \left( w_{2,1},w_{2,2},w_{2,3} \right)\) is the ideal they generate -- \({\overline{w}}_{2}\) is simply the vector whose components those three polynomials are (\cref{:rem-red-tau-omitted} explains why it is \({\overline{w}}_{2}\), not \(\tau({\overline{w}}_{2}) = \text{red}\left( \left\lbrack A_{1},A_{2} \right\rbrack \right)\) itself, that plays this role). At \(k = 3\), the three components of \({\overline{w}}_{3} = \tau({\overline{w}}_{2}) \times {\overline{v}}_{3}\) are \[2v_{1,2}v_{2,1}v_{3,2} - 2v_{1,1}v_{2,2}v_{3,2} - v_{1,3}v_{2,2}v_{3,3} + v_{1,2}v_{2,3}v_{3,3},\] \[- 2v_{1,2}v_{2,1}v_{3,1} + 2v_{1,1}v_{2,2}v_{3,1} - v_{1,3}v_{2,1}v_{3,3} + v_{1,1}v_{2,3}v_{3,3},\] \[v_{1,3}v_{2,2}v_{3,1} - v_{1,2}v_{2,3}v_{3,1} + v_{1,3}v_{2,1}v_{3,2} - v_{1,1}v_{2,3}v_{3,2},\] and \(I_{3} ≔ \left( w_{3,1},w_{3,2},w_{3,3} \right)\) likewise. At \(k = 4\), \({\overline{w}}_{4} = \tau({\overline{w}}_{3}) \times {\overline{v}}_{4}\): the same pattern continues, now quartic with up to 8 terms per component -- e.g. the first component of \({\overline{w}}_{4}\) is \[- 2v_{1,3}v_{2,2}v_{3,1}v_{4,2} + 2v_{1,2}v_{2,3}v_{3,1}v_{4,2} - 2v_{1,3}v_{2,1}v_{3,2}v_{4,2} + 2v_{1,1}v_{2,3}v_{3,2}v_{4,2}\] \[+ 2v_{1,2}v_{2,1}v_{3,2}v_{4,3} - 2v_{1,1}v_{2,2}v_{3,2}v_{4,3} - v_{1,3}v_{2,2}v_{3,3}v_{4,3} + v_{1,2}v_{2,3}v_{3,3}v_{4,3},\] the other two of the same shape (omitted here for space; all three components of \({\overline{w}}_{4}\), and \(I_{4} ≔ \left( w_{4,1},w_{4,2},w_{4,3} \right)\), computed and verified in code/verify\_determinantal\_structure\_all\_k.sage).

\end{example}

Stated over \(\mathbb{Q}\) since that is what the computational checks below actually run over, but nothing in the construction uses \(\mathbb{Q}\) specifically: every result in this article is proven for \(R_{k} ≔ \mathbb{F}\left\lbrack {\overline{v}}_{1},\ldots,{\overline{v}}_{k} \right\rbrack\) over an \textbf{arbitrary} field \(\mathbb{F}\) of characteristic \(\neq 2\) (needed only for \(\tau\)`s invertibility), by the identical argument with \(\mathbb{Q}\) replaced by \(\mathbb{F}\) throughout.

\section{The determinantal structure and the generator count}

\begin{theorem}{determinantal structure}{auto1} For \(k \geq 3\), \(I_{k}\) is exactly the ideal of \(2 \times 2\) minors of the \(2 \times 3\) matrix \[\begin{pmatrix}
\tau({\overline{w}}_{k - 1}) \\
{\overline{v}}_{k}
\end{pmatrix},\] i.e. \(V\left( I_{k} \right)\) is the preimage, under \(\left( {\overline{v}}_{1},\ldots,{\overline{v}}_{k} \right) \mapsto \begin{pmatrix}
\tau({\overline{w}}_{k - 1}\left( {\overline{v}}_{1},\ldots,{\overline{v}}_{k - 1} \right)) \\
{\overline{v}}_{k}
\end{pmatrix}\), of the classical rank-\(\leq 1\) determinantal variety of \(2 \times 3\) matrices. (At \(k = 2\) the same statement holds with \(\tau({\overline{w}}_{k - 1})\) replaced by \({\overline{v}}_{1}\) itself, matching \cref{:def-wk}`s base case \({\overline{w}}_{2} = {\overline{v}}_{1} \times {\overline{v}}_{2}\) directly.)

\end{theorem} \protect\phantomsection\label{thm-determinantal}{}

\begin{proof}[Proof of \cref{:thm-determinantal}] Immediate from the fact that the three components of a cross product \(a \times b\) are (up to sign) exactly the \(2 \times 2\) minors of the \(2 \times 3\) matrix with rows \(a,b\), combined with \cref{:def-wk}`s recursion -- and checked independently: comparing an independent implementation of the \(\tau\)/cross-product recursion against the ideal actually produced by the chained bilinear formula gives the \textbf{exact same ideal} (Gröbner-basis containment both ways) at \(k = 2,\ldots,6\) (code/verify\_determinantal\_structure\_all\_k.sage); the underlying polynomial identity \(\text{red}\left( \lbrack A,B\rbrack \right) = \tau(\text{red}(A) \times \text{ red}(B))\) makes this a two-line induction valid for \textbf{every} \(k\), not merely a pattern checked at small \(k\) (code/verify\_tau\_cross\_identity.sage).

\end{proof}

This theorem is the engine behind the geometric results of §5: \(R_{k}/I_{k}\) is not merely \textbf{determinantal-looking}, it is literally an iterated extension of the coordinate ring of a classical determinantal variety.

\begin{example}{the \(2 \times 3\) matrix of \cref{:thm-determinantal}, for \(k = 2,3,4\)}{lem-nondeg} At \(k = 2\): \[\begin{pmatrix}
{\overline{v}}_{1} \\
{\overline{v}}_{2}
\end{pmatrix} = \begin{pmatrix}
v_{1,1} & v_{1,2} & v_{1,3} \\
v_{2,1} & v_{2,2} & v_{2,3}
\end{pmatrix}.\] At \(k = 3\), the first row is \(\tau({\overline{w}}_{2}) = \left( w_{2,2},w_{2,1},2w_{2,3} \right)\) where \({\overline{w}}_{2} = {\overline{v}}_{1} \times {\overline{v}}_{2}\) is as in the previous example, giving \[\begin{pmatrix}
\tau({\overline{w}}_{2}) \\
{\overline{v}}_{3}
\end{pmatrix} = \begin{pmatrix}
v_{1,3}v_{2,1} - v_{1,1}v_{2,3} & - v_{1,3}v_{2,2} + v_{1,2}v_{2,3} & - 2v_{1,2}v_{2,1} + 2v_{1,1}v_{2,2} \\
v_{3,1} & v_{3,2} & v_{3,3}
\end{pmatrix}.\] At \(k = 4\), the first row is \(\tau({\overline{w}}_{3})\), entries cubic with up to 4 terms: \[\begin{pmatrix}
 - 2v_{1,2}v_{2,1}v_{3,1} + 2v_{1,1}v_{2,2}v_{3,1} - v_{1,3}v_{2,1}v_{3,3} + v_{1,1}v_{2,3}v_{3,3} & \quad\cdots & \quad\cdots \\
v_{4,1} & v_{4,2} & v_{4,3}
\end{pmatrix}\] (the other two entries of the first row, of the same shape, omitted for width -- all displayed in full in code/verify\_determinantal\_structure\_all\_k.sage). The row degree grows by one at each step (\(1,2,3,\ldots\)), while the second row is always the fresh, free vector \({\overline{v}}_{k}\) -- exactly the ``one fixed complicated row, one free row'' shape \cref{:prop-detideal} needs.

\end{example}

\begin{lemma}{non-degeneracy}{auto2} For every \(j \geq 2\), the polynomial map \({\overline{w}}_{j}:\mathbb{A}^{3j} \rightarrow \mathbb{A}^{3}\) is not identically zero. Hence \(U_{j} ≔ \left\{ p \in \mathbb{A}^{3j}:{\overline{w}}_{j}(p) \neq 0 \right\}\) is a nonempty Zariski-open subset.

\end{lemma} \protect\phantomsection\label{lem-nondeg}{}

\begin{proof}[Proof of \cref{:lem-nondeg}] By induction: \({\overline{w}}_{2}\left( e_{1},e_{2} \right) = e_{1} \times e_{2} = e_{3} \neq 0\); given \(c ≔ {\overline{w}}_{j - 1}(p) \neq 0\), at least one standard basis vector is not proportional to \(\tau(c)\), and a cross product vanishes only on linearly dependent pairs, so \({\overline{w}}_{j\left( p,e_{i} \right)} = \tau(c) \times e_{i} \neq 0\) for a suitable \(e_{i}\).

\end{proof}

The generator count of \(I_{k}\) needs nothing beyond the definition and non-degeneracy -- not codimension, not the determinantal-structure theorem above, not any external citation:

\begin{lemma}{independent same-degree generators are minimal}{lem-mingens-generic} If a homogeneous ideal in a graded ring is generated by finitely many forms of the same degree that are linearly independent over the base field, that generating set is minimal.

\end{lemma} \protect\phantomsection\label{lem-mingens-generic}{}

\begin{proof}[Proof of \cref{:lem-mingens-generic}] A polynomial-coefficient combination reproducing one of the generators from the others, while staying in the same degree, can only use degree-0 (constant) coefficients -- so a redundant generator would violate linear independence directly.

\end{proof}

\begin{lemma}{the three generators of \(I_{k}\) are independent}{lem-mingens-indep} The three components of \({\overline{w}}_{k}\) are linearly independent over \(\mathbb{F}\), for every \(k \geq 2\).

\end{lemma} \protect\phantomsection\label{lem-mingens-indep}{}

\begin{proof}[Proof of \cref{:lem-mingens-indep}] Checked directly (coefficient-matrix rank 3) at \(k = 2,\ldots,6\). Structurally: a nontrivial constant linear relation among the components of \({\overline{w}}_{k} = \tau({\overline{w}}_{k - 1}) \times {\overline{v}}_{k}\) would force \(\tau({\overline{w}}_{k - 1})\) itself to be a constant vector, which it is not for \(k \geq 3\) (\cref{:lem-nondeg}`s non-degeneracy, \(\tau\) invertible) nor for \(k = 2\) (\({\overline{v}}_{1}\) is a free variable block, not a constant).

\end{proof}

\begin{theorem}{{}\(I_{k}\) has exactly 3 minimal generators}{thm-mingens} For every \(k \geq 2\): \(I_{k}\) has exactly 3 minimal generators.

\end{theorem} \protect\phantomsection\label{thm-mingens}{}

\begin{proof}[Proof of \cref{:thm-mingens}] Immediate from \cref{:lem-mingens-generic} applied to \cref{:lem-mingens-indep}`s independence, since \(I_{k}\) is generated by the three (same-degree-\(k\)-homogeneous) components of \({\overline{w}}_{k}\) by definition.

\end{proof}

\section{\texorpdfstring{The geometric structure of \(V\left( I_{k} \right)\)}{The geometric structure of V\textbackslash left( I\_\{k\} \textbackslash right)}}

\begin{lemma}{continuous images and closures preserve irreducibility}{lem-top-irred} Let \(f:X \rightarrow Y\) be a continuous map of topological spaces with \(X\) irreducible. Then \(f(X)\) (with the subspace topology) is irreducible, and hence so is \(\overline{f(X)}\).

\end{lemma} \protect\phantomsection\label{lem-top-irred}{}

\begin{proof}[Proof of \cref{:lem-top-irred}] \textbf{Image.} Suppose \(f(X) = C_{1} \cup C_{2}\) with \(C_{1},C_{2}\) closed in \(f(X)\). By continuity, \(f^{- 1}\left( C_{1} \right)\) and \(f^{- 1}\left( C_{2} \right)\) are closed in \(X\), and \(X = f^{- 1}\left( C_{1} \right) \cup f^{- 1}\left( C_{2} \right)\). Since \(X\) is irreducible, one of the two equals \(X\) -- say \(f^{- 1}\left( C_{1} \right) = X\) -- so \(f(X) \subseteq C_{1}\), i.e. \(C_{1} = f(X)\), not proper. Hence \(f(X)\) is irreducible. \textbf{Closure.} If \(A \subseteq Y\) is irreducible and \(\overline{A} = C_{1} \cup C_{2}\) with \(C_{1},C_{2}\) closed in \(Y\), then \(A = \left( A \cap C_{1} \right) \cup \left( A \cap C_{2} \right)\); irreducibility of \(A\) forces, say, \(A \subseteq C_{1}\); since \(C_{1}\) is closed, \(\overline{A} \subseteq C_{1}\), and combined with \(C_{1} \subseteq \overline{A}\), \(C_{1} = \overline{A}\), not proper. Apply both parts with \(A ≔ f(X)\).

\end{proof}

This is purely topological -- continuity of \(f\) is the only hypothesis used, nothing about \(f\) being a polynomial map or morphism of varieties Youcis2014 (working-notes/ topological-irreducibility-of-image.org has the fuller discussion, including why this observation matters for a possible future Lean formalization of this article).

\begin{definition}{the strata \(Y_{1}\), \(Y_{2}\)}{def-y1y2} For \(k \geq 3\), write \(\pi\) for the projection \textbf{forgetting \({\overline{v}}_{k}\)}, \[\pi:\mathbb{A}^{3k} \rightarrow \mathbb{A}^{3(k - 1)},\quad\pi({\overline{v}}_{1},\ldots,{\overline{v}}_{k}) ≔ \left( {\overline{v}}_{1},\ldots,{\overline{v}}_{k - 1} \right),\] and \(\Phi\) for the map \[\Phi:\mathbb{A}^{3(k - 1)} \times \mathbb{A}^{1} \rightarrow \mathbb{A}^{3k},\quad\Phi(p,\lambda) ≔ \left. \begin{array}{r}
(p, \\
\lambda\tau({\overline{w}}_{k - 1}(p)))
\end{array} \right..\] Define two subsets of \(\mathbb{A}^{3k}\): \[Y_{1} ≔ V\left( I_{k - 1} \right) \times \mathbb{A}^{3},\quad Y_{2} ≔ \overline{\Phi(\mathbb{A}^{3k - 2})}.\] \(Y_{1}\) is the locus where the previous bracket already vanishes, with \({\overline{v}}_{k}\) completely free; \(Y_{2}\) is the closure of the locus swept out as \({\overline{v}}_{k}\) ranges over scalar multiples of \(\tau({\overline{w}}_{k - 1}(p))\), over all \(p \in \mathbb{A}^{3(k - 1)}\). At \(k = 2\) there is no ``previous bracket'' to chain from, so neither \(Y_{1}\) nor \(\Phi\) is defined; by convention \(Y_{2} ≔ V\left( I_{2} \right)\) directly there, matching \cref{:thm-determinantal}`s own base-case treatment.

\end{definition} \protect\phantomsection\label{def-y1y2}{}

\begin{figure}
\centering
\(\mathbb{A}^{3(k - 1)} \times \mathbb{A}^{1}\)

\(\Phi\) ↓ { \(\mathbb{A}^{3k}\) { \(Y_{1}\) } { \(Y_{2}\) } }
\caption{Both strata of \(V\left( I_{k} \right) = Y_{1} \cup Y_{2}\) inside the same ambient \(\mathbb{A}^{3k}\): \(Y_{1} ≔ V\left( I_{k - 1} \right) \times \mathbb{A}^{3}\) (identifying \(\mathbb{A}^{3(k - 1)} \times \mathbb{A}^{3} \cong \mathbb{A}^{3k}\)) and \(Y_{2} ≔ \overline{\Phi(\mathbb{A}^{3k - 2})}\), \(\Phi\)`s image (\(\Phi\) shown mapping in from above, into \(Y_{2}\) specifically -- \(Y_{1}\) is not part of \(\Phi\)`s image). The two overlap (not disjoint, Propositions 17-18 rule out either containing the other): \(\Phi(p,0) = (p,0)\) for \textbf{every} \(p\), so the whole zero-section \(\{(p,0)\}\) lies in \(Y_{2}\), and its intersection with \(Y_{1}\) is \(V\left( I_{k - 1} \right) \times \{ 0\}\), of dimension \(3k - 5\) -- properly smaller than either stratum's own dimension \(3k - 2\).}
\end{figure}

\protect\phantomsection\label{fig-phi-y2}{}

\begin{lemma}{{}\(Y_{2}\) is irreducible}{lem-y2-irred} For \(k \geq 3\): \(Y_{2}\) (\cref{:def-y1y2}) is irreducible, and in fact its defining ideal (the graph ideal of \(\Phi\) with \(\lambda\) eliminated) is prime.

\end{lemma} \protect\phantomsection\label{lem-y2-irred}{}

\begin{proof}[Proof of \cref{:lem-y2-irred}] \(\Phi\) is a polynomial map, hence continuous for the Zariski topology (preimages of Zariski-closed sets under a polynomial map are again Zariski-closed, by substitution), and its domain \(\mathbb{A}^{3k - 2}\) is irreducible (the polynomial ring over \(\mathbb{F}\) is a domain, so the zero ideal is prime). \cref{:lem-top-irred} gives \(Y_{2}\) irreducible directly. Primality of the elimination ideal follows the same way: it is the kernel of a map into a domain (\(\mathbb{F}\left\lbrack {\overline{v}}_{1},\ldots,{\overline{v}}_{k} \right\rbrack/( \cdot ) \cong \mathbb{F}\left\lbrack {\overline{v}}_{1},\ldots,{\overline{v}}_{k - 1},\lambda \right\rbrack\)), hence prime -- this holds for every \(k\) by this argument alone, additionally corroborated computationally at \(k = 4\) by computing the elimination ideal directly and confirming primality in Sage, as an independent check rather than a substitute for the argument.

\end{proof}

\begin{example}{the defining ideals of \(Y_{1}\) \textbf{and} \(Y_{2}\), for \(k = 2,3,4\)}{rem-closure-suffices} Writing \({\overline{v}}_{i} = \left( v_{i,1},v_{i,2},v_{i,3} \right)\) for the three coordinates of block \({\overline{v}}_{i}\): at \(k = 2\), \(Y_{2} ≔ V\left( I_{2} \right)\) (the base-case convention above), so its ideal \textbf{is} \(I_{2}\), generated by \[v_{1,2}v_{2,3} - v_{1,3}v_{2,2},\quad v_{1,3}v_{2,1} - v_{1,1}v_{2,3},\quad v_{1,1}v_{2,2} - v_{1,2}v_{2,1}.\] (\(Y_{1}\) is not defined at \(k = 2\), \cref{:def-y1y2}.) For \(k \geq 3\), \(Y_{1}\)`s defining ideal is simply \(I_{k - 1}\) itself, viewed inside \(R_{k}\) by leaving \({\overline{v}}_{k}\) free -- at \(k = 3\), exactly the three generators displayed above for \(I_{2}\), now inside \(R_{3}\) with \({\overline{v}}_{3}\) not appearing at all; at \(k = 4\), exactly \(I_{3}\)`s three generators (displayed below in this same example) inside \(R_{4}\) with \({\overline{v}}_{4}\) free. Unlike \(Y_{2}\), no elimination is needed to compute \(Y_{1}\)`s ideal -- it is already \(I_{k - 1}\)`s own generating set, unchanged. At \(k = 3\), \(Y_{2}\)`s elimination ideal (computed directly, not inferred) is generated by \[2v_{2,2}v_{3,1} + 2v_{2,1}v_{3,2} + v_{2,3}v_{3,3},\quad 2v_{1,2}v_{3,1} + 2v_{1,1}v_{3,2} + v_{1,3}v_{3,3},\] \[2v_{1,2}v_{2,1}v_{3,2} - 2v_{1,1}v_{2,2}v_{3,2} - v_{1,3}v_{2,2}v_{3,3} + v_{1,2}v_{2,3}v_{3,3}\] -- notably, two of the three generators are \textbf{linear} in \({\overline{v}}_{3}\) (exactly the two components of the ``proportionality'' condition \({\overline{v}}_{3} \propto \tau({\overline{w}}_{2})\)), and only the third, not involving \({\overline{v}}_{3}\) at all, carries the full complexity of \(\tau({\overline{w}}_{2})\) itself. \textbf{This ideal is a proper factor of \(I_{3}\), not equal to it} -- \(I_{3}\) is the intersection of this ideal with a second prime (the ``\(Y_{1}\)'' piece, where the \(k = 2\) bracket already vanishes); see the component-count example after \cref{:thm-componentcount} below, where this exact ideal reappears as one of \(I_{3}\)`s two primary components. At \(k = 4\), by the same pattern, two of the three generators are linear in \({\overline{v}}_{4}\) and the third is not, now degree 4: \[2v_{3,2}v_{4,1} + 2v_{3,1}v_{4,2} + v_{3,3}v_{4,3},\] \[v_{1,3}v_{2,2}v_{4,1} - v_{1,2}v_{2,3}v_{4,1} - v_{1,3}v_{2,1}v_{4,2} + v_{1,1}v_{2,3}v_{4,2} + v_{1,2}v_{2,1}v_{4,3} - v_{1,1}v_{2,2}v_{4,3},\] \[2v_{1,3}v_{2,2}v_{3,1}v_{4,2} - 2v_{1,2}v_{2,3}v_{3,1}v_{4,2} + 2v_{1,3}v_{2,1}v_{3,2}v_{4,2} - 2v_{1,1}v_{2,3}v_{3,2}v_{4,2}\] \[- 2v_{1,2}v_{2,1}v_{3,2}v_{4,3} + 2v_{1,1}v_{2,2}v_{3,2}v_{4,3} + v_{1,3}v_{2,2}v_{3,3}v_{4,3} - v_{1,2}v_{2,3}v_{3,3}v_{4,3}.\] All computed directly (elimination of \(\lambda\) from the graph ideal of \(\Phi\), in Sage), not just asserted from the pattern.

\end{example}

\begin{remark}{the closure, not the raw image, is what is needed}{auto3} A polynomial map's image need not itself be closed -- the standard example is \(\mathbb{A}^{2} \rightarrow \mathbb{A}^{2}\), \((x,y) \mapsto (x,xy)\), whose image is not even locally closed, though it is \textbf{constructible} (a finite union of locally closed sets) by Chevalley's theorem, which has no purely topological analogue for general continuous maps Youcis2014. This never causes a problem for \(Y_{2}\), because \(Y_{2}\) is \textbf{defined} as \(\overline{\Phi( \cdot )}\) from the outset, and the decomposition used in \cref{:thm-codim}`s proof below, \(V\left( I_{k} \right) = Y_{1} \cup Y_{2}\), is established by a direct case split rather than by taking closures and hoping: a point \(\left( p,{\overline{v}}_{k} \right) \in V\left( I_{k} \right)\) either has \(\tau({\overline{w}}_{k - 1}(p)) = 0\), putting it in \(Y_{1}\) directly, or has \({\overline{v}}_{k} = \lambda\tau({\overline{w}}_{k - 1}(p))\) for some \(\lambda\), so \(\left( p,{\overline{v}}_{k} \right) = \Phi(p,\lambda)\) is \textbf{literally} a value of \(\Phi\) -- not merely a limit point -- hence in \(Y_{2}\) with no appeal to closure at all. The reverse containment \(Y_{1} \cup Y_{2} \subseteq V\left( I_{k} \right)\) needs only that \(\Phi( \cdot ) \subseteq V\left( I_{k} \right)\) (a cross product of a vector with itself vanishes: \(\tau({\overline{w}}_{k - 1}) \times \left( \lambda\tau({\overline{w}}_{k - 1}) \right) = 0\) identically) together with \(V\left( I_{k} \right)\) already being closed, automatic since it is a vanishing locus by definition. So the theorem never needs ``the image of \(\Phi\) is closed'' -- only that it is contained in the already-closed \(V\left( I_{k} \right)\), which is immediate.

\end{remark} \protect\phantomsection\label{rem-closure-suffices}{}

\begin{lemma}{dimension of \(Y_{2}\), and where it lives}{lem-y2-dim} With \(Y_{2}\) as in \cref{:lem-y2-irred} and \(U_{k - 1}\) as in \cref{:lem-nondeg}: \(\dim Y_{2} = 3k - 2\), and \(Y_{2}\) meets \(\pi^{- 1}\left( U_{k - 1} \right)\) nonemptily, where \(\pi\) is the projection of \cref{:def-y1y2}, forgetting \({\overline{v}}_{k}\).

\end{lemma} \protect\phantomsection\label{lem-y2-dim}{}

\begin{proof}[Proof of \cref{:lem-y2-dim}] \(\pi:\mathbb{A}^{3k} \rightarrow \mathbb{A}^{3(k - 1)}\) discards the last block: \(\pi({\overline{v}}_{1},\ldots,{\overline{v}}_{k}) = \left( {\overline{v}}_{1},\ldots,{\overline{v}}_{k - 1} \right)\). Hence \(\pi^{- 1}\left( U_{k - 1} \right) = U_{k - 1} \times \mathbb{A}^{3}\), the set of points whose \textbf{first} \(k - 1\) blocks lie in \(U_{k - 1}\) and whose \({\overline{v}}_{k}\) is completely unconstrained. Over \(U_{k - 1}\), \(\Phi\) is injective. Both \(p\) and \(\lambda\) can be recovered from \(\Phi(p,\lambda)\): \(p\) directly, and \(\lambda\) by dividing any nonzero coordinate of \({\overline{v}}_{k} = \lambda\tau({\overline{w}}_{k - 1}(p))\) by the corresponding coordinate of \(\tau({\overline{w}}_{k - 1}(p))\); this division is well-defined since \(\tau({\overline{w}}_{k - 1}(p)) \neq 0\) on \(U_{k - 1}\). Injectivity gives \(\dim Y_{2} = \dim\left( U_{k - 1} \times \mathbb{A}^{1} \right) = 3(k - 1) + 1 = 3k - 2\). Since \(U_{k - 1}\) is nonempty by \cref{:lem-nondeg}, \(\Phi(U_{k - 1} \times \left\{ 1 \right\})\) is a nonempty subset of \(Y_{2} \cap \pi^{- 1}\left( U_{k - 1} \right)\).

\end{proof}

\begin{theorem}{{}\(I_{k}\) has codimension 2}{thm-codim} For every \(k \geq 2\): \(\dim V\left( I_{k} \right) = 3k - 2\), i.e. \(\text{codim}\left( I_{k} \right) = 2\).

\end{theorem} \protect\phantomsection\label{thm-codim}{}

\begin{proof}[Proof of \cref{:thm-codim}] Induction on \(k\). \textbf{Base case \(k = 2\)}: \(V\left( I_{2} \right)\) is the rank-\(\leq 1\) locus of a \(2 \times 3\) matrix with independent-indeterminate entries, of the standard dimension \(4 = 3 \cdot 2 - 2\) for the generic determinantal variety of this shape MotzkinTaussky1955. \textbf{Inductive step}: by \cref{:thm-determinantal}, and since \(a \times b = 0\) if and only if \(a,b\) are linearly dependent, \(V\left( I_{k} \right) = Y_{1} \cup Y_{2}\) (\cref{:def-y1y2}). The locus \(\tau({\overline{w}}_{k - 1}) = 0\) gives \(Y_{1}\), with \({\overline{v}}_{k}\) unconstrained; the locus \(\tau({\overline{w}}_{k - 1}) \neq 0\) with \({\overline{v}}_{k}\) proportional to it gives a subset of \(Y_{2}\). This is the direct case split of \cref{:rem-closure-suffices}, not merely a closure-level approximation. By the inductive hypothesis, \(\dim V\left( I_{k - 1} \right) = 3(k - 1) - 2\), so \(\dim Y_{1} = 3(k - 1) - 2 + 3 = 3k - 2\). Also \(\dim Y_{2} = 3k - 2\), by \cref{:lem-y2-dim}. Hence \(\dim V\left( I_{k} \right) = \dim(Y_{1} \cup Y_{2}) = 3k - 2\). This dimension statement alone needs no containment argument; that is needed only for the sharper component-count result below.

\end{proof}

\begin{example}{dimension and codimension, for \(k = 2,3,4\)}{prop-y2-not-in-y1} Values of \(\dim R_{k}\) (the ambient dimension, \(= 3k\)), \(\dim(R_{k}/I_{k})\) (the dimension of the variety \(V\left( I_{k} \right)\) -- \textbf{not} ``\(\dim I_{k}\)'', an ideal is not a variety and does not itself have a dimension in this sense), and \(\text{codim}\left( I_{k} \right) ≔ \dim R_{k} - \dim(R_{k}/I_{k})\):

\begin{longtable}[]{@{}cccc@{}}
\toprule\noalign{}
\endhead
\bottomrule\noalign{}
\endlastfoot
\(k\) & \(\dim R_{k}\) & \(\dim(R_{k}/I_{k})\) & \(\text{codim}\left( I_{k} \right)\) \\
2 & 6 & 4 & 2 \\
3 & 9 & 7 & 2 \\
4 & 12 & 10 & 2 \\
\end{longtable}

All three rows follow directly from \cref{:thm-codim}`s formula \(\dim(R_{k}/I_{k}) = 3k - 2\); also confirmed by direct Gröbner-basis dimension computation in Sage at each \(k\), as an independent check.

\end{example}

\begin{proposition}{{}\(Y_{2} \nsubseteq Y_{1}\)}{auto4} With \(Y_{1},Y_{2}\) as in the proof of \cref{:thm-codim}: \(Y_{2} \nsubseteq Y_{1}\).

\end{proposition} \protect\phantomsection\label{prop-y2-not-in-y1}{}

\begin{proof}[Proof of \cref{:prop-y2-not-in-y1}] Pick \(p \in U_{k - 1}\) (\cref{:lem-nondeg}); then \(\Phi(p,0) = (p,0) \in Y_{2}\) but \({\overline{w}}_{k - 1}(p) \neq 0\) puts \((p,0) \notin Y_{1}\).

\end{proof}

\begin{proposition}{no component of \(Y_{1}\) lies in \(Y_{2}\)}{prop-y1-not-in-y2} No irreducible component \(Z\) of \(V\left( I_{k - 1} \right)\) (each of dimension \(3(k - 1) - 2\), by \cref{:thm-codim} applied at \(k - 1\)) has \(Z \times \mathbb{A}^{3} \subseteq Y_{2}\).

\end{proposition} \protect\phantomsection\label{prop-y1-not-in-y2}{}

\begin{proof}[Proof of \cref{:prop-y1-not-in-y2}] Both \(Z \times \mathbb{A}^{3}\) and \(Y_{2}\) are irreducible: \(Y_{2}\) by \cref{:lem-y2-irred}, and \(Z \times \mathbb{A}^{3}\) because a product of irreducible varieties is irreducible (\(Z\) is irreducible by hypothesis, and \(\mathbb{A}^{3}\) trivially is). Both have the \textbf{same} dimension \(3k - 2\), by \cref{:thm-codim} and \cref{:lem-y2-dim}. Containment would therefore force equality. This is refuted: \(Z \times \mathbb{A}^{3}\) is disjoint from \(\pi^{- 1}\left( U_{k - 1} \right)\), since \({\overline{w}}_{k - 1}\) vanishes identically on \(Z \subset V\left( I_{k - 1} \right)\); but \(Y_{2}\) meets \(\pi^{- 1}\left( U_{k - 1} \right)\), by \cref{:lem-y2-dim}.

\end{proof}

\begin{theorem}{component count}{thm-componentcount} For every \(k \geq 2\): \(V\left( I_{k} \right)\) has exactly \(k - 1\) irreducible components, each of dimension \(3k - 2\). The component indexed \(j = 2,\ldots,k\) involves exactly the first \(j\) coordinate blocks and nothing later.

\end{theorem} \protect\phantomsection\label{thm-componentcount}{}

\begin{proof}[Proof of \cref{:thm-componentcount}] Induction on \(k\). \textbf{Base case \(k = 2\)}: the rank-\(\leq 1\) locus of a generic \(2 \times 3\) matrix is irreducible -- the affine cone over the Segre embedding \(\mathbb{P}^{1} \times \mathbb{P}^{2} \hookrightarrow \mathbb{P}^{5}\), equivalently the Motzkin-Taussky MotzkinTaussky1955 irreducibility of the \(2 \times 2\) commuting variety -- so exactly 1 component, matching \(k - 1 = 1\). \textbf{Step}: \(V\left( I_{k} \right) = Y_{1} \cup Y_{2}\) as in \cref{:thm-codim}`s proof; by the inductive hypothesis \(Y_{1}\) has \(k - 2\) irreducible components, each of dimension \(3k - 2\) (\cref{:thm-codim}), none contained in another; \cref{:prop-y2-not-in-y1} and \cref{:prop-y1-not-in-y2} rule out containment between \(Y_{2}\) and any of them or vice versa. A finite union of irreducible closed sets with no containments is irredundant, giving exactly \((k - 2) + 1 = k - 1\) components.

\end{proof}

This gives \(I_{k}\)`s minimal primes and their count, but not yet that \(I_{k}\) is radical (so not yet that these are literally the components of a \textbf{primary} decomposition of \(I_{k}\) itself, as opposed to of \(\sqrt{I_{k}}\)) -- that is closed in §6, after Cohen-Macaulayness is available.

\begin{definition}{the components \(P_{j}\)}{def-primecomponents} For \(k \geq 2\) and \(j = 2,\ldots,k\), write \(P_{j}\) for the vanishing ideal of the component of \(V\left( I_{k} \right)\) indexed \(j\) in \cref{:thm-componentcount} -- the one involving exactly the first \(j\) coordinate blocks. Explicitly, recursively in \(j\):

\begin{itemize}
\item
  \textbf{Base case, \(j = k\)}: \(P_{k}\) is \(Y_{2}\)`s own defining ideal at parameter \(k\) (\cref{:lem-y2-irred}).
\item
  \textbf{Recursive case, \(j < k\)}: \(P_{j}\) is that same ideal computed at parameter \(j\) instead -- i.e. \(Y_{2}\)`s defining ideal when the construction is run with \(k = j\) -- viewed inside \(R_{k}\) by leaving \({\overline{v}}_{j + 1},\ldots,{\overline{v}}_{k}\) free.
\item
  \textbf{Consequence}: \(P_{j}\) does not depend on which \(k \geq j\) it is viewed inside of -- it is the same ideal, just embedded in a ring with more or fewer free variables.
\end{itemize}

\end{definition} \protect\phantomsection\label{def-primecomponents}{}

\begin{example}{a suitable point of \(Y_{2}\), concretely, at \(k = 3\)}{prop-rational} Rather than only the defining ideal (Example 24 below), a single explicit point: take \({\overline{v}}_{1} = (1,0,0)\), \({\overline{v}}_{2} = (0,1,0)\) (so \({\overline{v}}_{1} \in U_{2}\), \cref{:lem-nondeg}), giving \({\overline{w}}_{2} = {\overline{v}}_{1} \times {\overline{v}}_{2} = (0,0,1)\) and \(\tau({\overline{w}}_{2}) = (0,0,2)\). Taking \(\lambda = 1\) in \cref{:def-y1y2}`s \(\Phi\), i.e. \({\overline{v}}_{3} ≔ \tau({\overline{w}}_{2}) = (0,0,2)\), gives the point \[\left( {\overline{v}}_{1},{\overline{v}}_{2},{\overline{v}}_{3} \right) = \left( (1,0,0),(0,1,0),(0,0,2) \right) \in Y_{2} \subset \mathbb{A}^{9.}\] Directly: \({\overline{w}}_{3} = \tau({\overline{w}}_{2}) \times {\overline{v}}_{3} = (0,0,2) \times (0,0,2) = 0\) (a vector crossed with itself), confirming this point lies in \(V\left( I_{3} \right)\), as any point of the form \(\Phi(p,\lambda)\) must (\cref{:def-y1y2}). This is also, concretely, a point of the largest-dimension inherited case: it is \textbf{not} a rank-2 witness in \cref{:lem-witness-new}`s sense (that witness uses \(\lambda = 0\), not \(\lambda = 1\)) -- it illustrates \(Y_{2}\) itself, not the radicality argument's own witness points.

\end{example}

\begin{proposition}{rational parametrization of each component}{auto5} For \(k \geq 2\) and \(j = 2,\ldots,k\), the component of \(V\left( I_{k} \right)\) indexed \(j\) (\cref{:thm-componentcount}) is a \textbf{rational} variety: it is the closure of the image of the polynomial map \[\mathbb{A}^{3(j - 1)} \times \mathbb{A}^{1} \times \mathbb{A}^{3(k - j)} \rightarrow \mathbb{A}^{3k},\quad(p,\lambda,q) \mapsto \left( \Phi_{j}(p,\lambda),q \right),\] where \(\Phi_{j}\) is \(\Phi\) (\cref{:def-y1y2}) run at parameter \(j\) in \cref{:def-primecomponents}'' sense (at \(j = 2\), \(\Phi_{2}(p,\lambda) ≔ (p,\lambda p)\), matching \cref{:thm-determinantal}`s own base-case convention), and \(q\) fills the remaining free blocks \({\overline{v}}_{j + 1},\ldots,{\overline{v}}_{k}\). This map is moreover \textbf{birational} onto the component: it is injective over \(U_{j - 1} \times \mathbb{A}^{1} \times \mathbb{A}^{3(k - j)}\), by the same argument as \cref{:lem-y2-dim}, applied at parameter \(j\).

\end{proposition} \protect\phantomsection\label{prop-rational}{}

\begin{proof}[Proof of \cref{:prop-rational}] By \cref{:def-primecomponents}, \(P_{j}\)`s defining ideal is \(Y_{2}\)`s own defining ideal run at parameter \(j\) -- i.e. the component itself \textbf{is} \(Y_{2}\) at parameter \(j\), crossed with the free blocks \({\overline{v}}_{j + 1},\ldots,{\overline{v}}_{k}\). \cref{:lem-y2-dim}`s injectivity argument (over \(U_{j - 1}\), via \cref{:lem-nondeg} at level \(j - 1\)) applies verbatim at parameter \(j\), and appending the free blocks \(q\) unchanged does not affect injectivity.

\end{proof}

\begin{remark}{a reducible variety is not itself ``rational'' -- component-by-component is what \cref{:prop-rational} answers}{prop-detideal} \(Y_{1} ≔ V\left( I_{k - 1} \right) \times \mathbb{A}^{3}\) (\cref{:def-y1y2}) is \textbf{not} irreducible for \(k \geq 4\) -- it is a union of \(k - 2\) of the components above (\cref{:thm-componentcount}`s inductive step) -- so a single rational parametrization \textbf{of \(Y_{1}\) as one object} is not meaningful in the usual sense (rationality is a property of irreducible varieties). What \textbf{is} meaningful, and what \cref{:prop-rational} supplies, is a rational parametrization of each of \(Y_{1}\)`s irreducible pieces individually -- exactly the components inherited from \(V\left( I_{k - 1} \right)\).

\end{remark}

\begin{example}{the components, worked out at \(k = 3\)}{auto6} Uniformly by \(k\), since the components'' \textbf{count} and \textbf{block pattern} are already known (\cref{:thm-componentcount}, \cref{:def-primecomponents}) and only their explicit generators change:

\begin{itemize}
\item
  \textbf{At \(k = 2\)}: \(I_{2}\) is already prime (\cref{:thm-componentcount}`s base case), so \(V\left( I_{2} \right)\) has exactly \(k - 1 = 1\) component -- itself.
\item
  \textbf{At \(k = 3\)}: exactly \(k - 1 = 2\) components, both computed directly (Gröbner bases, not inferred from the pattern alone), matching \(I_{3} = P_{2} \cap P_{3}\):

  \begin{itemize}
  \item
    \(P_{2}\), the ``\(Y_{1}\)'' piece (block pattern \(\{ 1,2\}\), matching that the length-2 sub-bracket already vanishes with \({\overline{v}}_{3}\) free) -- generated by exactly \(I_{2}\)`s three generators from the first example, now inside \(R_{3}\) with \({\overline{v}}_{3}\) simply not appearing: \[v_{1,2}v_{2,3} - v_{1,3}v_{2,2},\quad v_{1,3}v_{2,1} - v_{1,1}v_{2,3},\quad v_{1,1}v_{2,2} - v_{1,2}v_{2,1}.\]
  \item
    \(P_{3}\), the ``\(Y_{2}\)'' piece (block pattern \(\{ 1,2,3\}\)) -- exactly the ideal already computed in the \(Y_{2}\)-defining-ideal example above: \[2v_{2,2}v_{3,1} + 2v_{2,1}v_{3,2} + v_{2,3}v_{3,3},\quad 2v_{1,2}v_{3,1} + 2v_{1,1}v_{3,2} + v_{1,3}v_{3,3},\] \[2v_{1,2}v_{2,1}v_{3,2} - 2v_{1,1}v_{2,2}v_{3,2} - v_{1,3}v_{2,2}v_{3,3} + v_{1,2}v_{2,3}v_{3,3}.\]
  \end{itemize}
\item
  \textbf{At \(k = 4\)}: exactly \(k - 1 = 3\) components, block patterns \(\{ 1,2\}\), \(\{ 1,2,3\}\), \(\{ 1,2,3,4\}\):

  \begin{itemize}
  \item
    The first two, \(P_{2}\) and \(P_{3}\), are inherited \textbf{verbatim} from \(k = 3\) above, with \({\overline{v}}_{4}\) simply left free -- not recomputed, per \cref{:def-primecomponents}'' own ``does not depend on which \(k\)'' clause.
  \item
    The third, \(P_{4}\), is the new ``\(Y_{2}\)'' stratum -- exactly the \(k = 4\) row of the earlier \(Y_{2}\)-defining-ideal example (generators omitted here to avoid repeating that display; same growth pattern, quartic with up to 8 terms per generator).
  \end{itemize}
\end{itemize}

\end{example}

\section{Cohen-Macaulayness}

\begin{proposition}{Eagon-Northcott / Hilbert-Burch, applied}{auto7} Let \(M\) be a \(2 \times 3\) matrix with entries in a polynomial ring \(R\) over a field (arbitrary entries, not necessarily independent indeterminates), and \(J\) the ideal of its \(2 \times 2\) minors. If \(J\) has the \textbf{expected} codimension \(3 - 2 + 1 = 2\), then \(J\) is a perfect ideal of grade 2 -- \(\text{pd}(R/J) = 2\), hence \(R/J\) is Cohen-Macaulay.

\end{proposition} \protect\phantomsection\label{prop-detideal}{}

\begin{proof}[Proof of \cref{:prop-detideal}] Bruns and Vetter BrunsVetter1988, Theorem (2.7), p. 13, gives exactly this for an \textbf{arbitrary} matrix, as needed since one row of our \(2 \times 3\) matrix (\(\tau({\overline{w}}_{k - 1})\)) is a fixed polynomial vector, not a row of independent indeterminates. (The original construction, Eagon and Northcott EagonNorthcott1962, is superseded here by Bruns-Vetter's more general treatment -- confirmed closed-access via Unpaywall, not used directly.)

\end{proof}

\begin{remark}{why the generator count above wasn't proven this way}{thm-cm} It is tempting to expect \cref{:prop-detideal} to also hand over the generator count for free via Bruns-Vetter's Theorem (16.36), p. 217 (``The Theorem of Hilbert-Burch'') -- but that theorem, as actually stated, gives only \(I = a \cdot I_{m}(f)\) for a nonzerodivisor \(a\), not a generator count on its own: at expected codimension, an \textbf{arbitrary} \(2 \times 3\) matrix's minors can be perfect and Cohen-Macaulay while one minor vanishes identically or is a scalar combination of the others (concretely: \(\begin{pmatrix}
x & y & 0 \\
0 & 0 & 1
\end{pmatrix}\) over \(\mathbb{F}\lbrack x,y\rbrack\) has minors \(0,x,y\), height 2 as expected, \(J = (x,y)\) perfect and Cohen-Macaulay with only \textbf{2} generators). This is exactly why \cref{:thm-mingens} was proven directly from linear independence instead of cited from this section's machinery.

\end{remark}

\begin{theorem}{{}\(R_{k}/I_{k}\) is Cohen-Macaulay}{auto8} For every \(k \geq 2\): \(R_{k}/I_{k}\) is Cohen-Macaulay.

\end{theorem} \protect\phantomsection\label{thm-cm}{}

\begin{proof}[Proof of \cref{:thm-cm}] By \cref{:thm-codim}, \(I_{k}\) has codimension 2, the expected value for the \(2 \times 3\) matrix of \cref{:thm-determinantal} presenting it. \cref{:prop-detideal} applies directly. (This uses \cref{:thm-codim} only -- not \cref{:thm-mingens}`s generator count, not \cref{:thm-componentcount}`s component count -- genuinely independent of both.)

\end{proof}

\section{Radicality}

Cohen-Macaulayness and the explicit component structure of §5 combine, via Serre's criterion, to give radicality -- which then closes both the generator-count (§4) / Cohen-Macaulay (§6) story (making \(I_{k}\) itself, not just \(\sqrt{I_{k}}\), the object with those properties) and upgrades §5's component count to a genuine primary decomposition.

\begin{lemma}{reduction via Serre}{lem-serre} Recall Serre's condition \(S_{n}\): a Noetherian ring \(A\) satisfies \(S_{n}\) if \(\text{depth}\left( A_{P} \right) \geq \min(n,\dim A_{P})\) for every prime \(P \subset A\) (Cohen-Macaulay means \(\text{depth}\left( A_{P} \right) = \dim A_{P}\) at \textbf{every} prime -- the extremal case, \(n = \infty\) in effect). Cohen-Macaulay rings satisfy \(S_{n}\) for every \(n\): \(\text{depth}\left( A_{P} \right) = \dim A_{P} \geq \min(n,\dim A_{P})\) trivially, whatever \(n\) is -- this needs nothing beyond the two definitions just given, standard for Serre's conditions in general (see e.g. Bruns and Vetter BrunsVetter1988 or Eisenbud's \emph{Commutative Algebra with a View Toward Algebraic Geometry}). In particular \(S_{1}\). So, given \cref{:thm-cm}, radicality of \(I_{k}\) is equivalent to \(R_{0}\): regularity of \(R_{k}/I_{k}\) at each of its minimal primes.

\end{lemma} \protect\phantomsection\label{lem-serre}{}

\begin{lemma}{Jacobian criterion: the rank bound}{lem-jacobian-rankbound} Let \(S = \mathbb{F}\left\lbrack x_{1},\ldots,x_{n} \right\rbrack\), \(I = \left( f_{1},\ldots,f_{m} \right)\), \(A = S/I\) equidimensional of dimension \(n - c\), and \(x \in V(I)\) an \(\mathbb{F}\)-rational point with Jacobian \(\text{Jac}(x) \in \mathbb{F}^{m \times n}\). Then \(\text{rank }\text{ Jac}(x) \leq c\).

\end{lemma} \protect\phantomsection\label{lem-jacobian-rankbound}{}

\begin{proof}[Proof of \cref{:lem-jacobian-rankbound}] \(\dim_{\mathbb{F}}\mathfrak{m}/\mathfrak{m}^{2} = n - \text{ rank }\text{ Jac}(x)\) always (linear algebra: the image of \(I\) in the cotangent space \(\mathfrak{n}/\mathfrak{n}^{2} \cong \mathbb{F}^{n}\) is exactly the span of the linear parts of generators of \(I\), of dimension \(\text{rank }\text{ Jac}(x)\)), and \(\dim_{\mathbb{F}}\mathfrak{m}/\mathfrak{m}^{2} \geq \dim A_{\mathfrak{m}}\) in any Noetherian local ring, so it suffices that \(\dim A_{\mathfrak{m}} = n - c\) -- which \cref{:lem-jacobian-regularity} supplies.

\end{proof}

\begin{lemma}{Jacobian criterion: regularity at rank \(= c\)}{lem-jacobian-regularity} With hypotheses as in \cref{:lem-jacobian-rankbound}: \(\dim A_{\mathfrak{m}} = n - c\). Consequently, if \(\text{rank }\text{ Jac}(x) = c\), then \(A_{\mathfrak{m}}\) is regular of dimension \(n - c\); in particular \(x\) lies on a \textbf{unique} component of \(V(I)\), whose local ring there is a field.

\end{lemma} \protect\phantomsection\label{lem-jacobian-regularity}{}

\begin{proof}[Proof of \cref{:lem-jacobian-regularity}] \(\dim A_{\mathfrak{m}} = \max\left\{ \dim(A/P)_{\mathfrak{m}}:P \in \text{ Min}(A),P \subseteq \mathfrak{m} \right\}\), and each \(A/P = S/\widetilde{P}\) is an affine domain over a field, for which \(\dim\left( S/\widetilde{P} \right)_{\mathfrak{M}} = \dim S/\widetilde{P}\) at \textbf{every} maximal ideal \(\mathfrak{M}\) -- affine domains over a field are catenary and equicodimensional (Eisenbud, \emph{Commutative Algebra}, Cor. 13.4; equivalently the Stacks Project, tag 00OS Stacks00OS -- a dimension formula for finitely generated field extensions, holding over an arbitrary field, not needing a perfect-field hypothesis by accident). By equidimensionality every such \(\dim S/\widetilde{P}\) equals \(n - c\). If \(\text{rank }\text{ Jac}(x) = c\) then \(\dim_{\mathbb{F}}\mathfrak{m}/\mathfrak{m}^{2} = n - c = \dim A_{\mathfrak{m}}\) by \cref{:lem-jacobian-rankbound}, the definition of regularity. A regular local ring is a domain, so \(\mathfrak{m}\) contains a unique minimal prime of \(A\), and the local ring at that prime is a further localization of a regular local ring at its own maximal ideal in dimension 0, hence a field.

\end{proof}

\begin{lemma}{the Jacobian's own recursion}{lem-jacrecursion} Write \(J_{k}\) for the \(3 \times 3k\) Jacobian of \({\overline{w}}_{k}\), \(\lbrack u\rbrack_{\times}\) for the skew matrix with \(\lbrack u\rbrack_{\times}v = u \times v\), and \(T\) for the matrix of \(\tau\). Then \(J_{2} = \left\lbrack - \left\lbrack {\overline{v}}_{2} \right\rbrack_{\times}|\left\lbrack {\overline{v}}_{1} \right\rbrack_{\times} \right\rbrack\), and for \(k \geq 3\), \[J_{k} = \left\lbrack - \left\lbrack {\overline{v}}_{k} \right\rbrack_{\times}TJ_{k - 1}|\left\lbrack \tau({\overline{w}}_{k - 1}) \right\rbrack_{\times} \right\rbrack.\]

\end{lemma} \protect\phantomsection\label{lem-jacrecursion}{}

\begin{proof}[Proof of \cref{:lem-jacrecursion}] For the second block, \({\overline{w}}_{k} = \tau({\overline{w}}_{k - 1}) \times {\overline{v}}_{k} = \left\lbrack \tau({\overline{w}}_{k - 1}) \right\rbrack_{\times}{\overline{v}}_{k}\) is linear in \({\overline{v}}_{k}\) with coefficient matrix free of \({\overline{v}}_{k}\). For the first, differentiating \({\overline{w}}_{k} = \tau({\overline{w}}_{k - 1}) \times {\overline{v}}_{k}\) with respect to any variable \(t\) among \({\overline{v}}_{1},\ldots,{\overline{v}}_{k - 1}\) gives \(\partial_{t}{\overline{w}}_{k} = \left( T\partial_{t}{\overline{w}}_{k - 1} \right) \times {\overline{v}}_{k} = - \left\lbrack {\overline{v}}_{k} \right\rbrack_{\times}T\partial_{t}{\overline{w}}_{k - 1}\), using \(a \times b = - \lbrack b\rbrack_{\times}a\); assembling columns gives the stated product.

\end{proof}

\begin{example}{{}\(J_{3}\), explicitly}{lem-witness-new} With variables ordered \(v_{1,1},v_{1,2},v_{1,3},v_{2,1},v_{2,2},v_{2,3},v_{3,1},v_{3,2},v_{3,3}\), the \(3 \times 9\) Jacobian of \({\overline{w}}_{3}\) (the three generators of \(I_{3}\) from the earlier example), column-blocked to match \cref{:lem-jacrecursion}`s formula \(J_{3} = \left\lbrack - \left\lbrack {\overline{v}}_{3} \right\rbrack_{\times}TJ_{2}|\left\lbrack \tau({\overline{w}}_{2}) \right\rbrack_{\times} \right\rbrack\) -- shown here as the two blocks separately, since side by side the full \(3 \times 9\) matrix does not fit the page width. First, \(- \left\lbrack {\overline{v}}_{3} \right\rbrack_{\times}TJ_{2}\) (\(3 \times 6\)): \[\begin{pmatrix}
 - 2v_{2,2}v_{3,2} & 2v_{2,1}v_{3,2} + v_{2,3}v_{3,3} & - v_{2,2}v_{3,3} & 2v_{1,2}v_{3,2} & - 2v_{1,1}v_{3,2} - v_{1,3}v_{3,3} & v_{1,2}v_{3,3} \\
2v_{2,2}v_{3,1} + v_{2,3}v_{3,3} & - 2v_{2,1}v_{3,1} & - v_{2,1}v_{3,3} & - 2v_{1,2}v_{3,1} - v_{1,3}v_{3,3} & 2v_{1,1}v_{3,1} & v_{1,1}v_{3,3} \\
 - v_{2,3}v_{3,2} & - v_{2,3}v_{3,1} & v_{2,2}v_{3,1} + v_{2,1}v_{3,2} & v_{1,3}v_{3,2} & v_{1,3}v_{3,1} & - v_{1,2}v_{3,1} - v_{1,1}v_{3,2}
\end{pmatrix}\] and then \(\left\lbrack \tau({\overline{w}}_{2}) \right\rbrack_{\times}\) (\(3 \times 3\), the skew matrix of \(\tau({\overline{w}}_{2})\)`s three coordinates from the earlier matrix example): \[\begin{pmatrix}
0 & 2v_{1,2}v_{2,1} - 2v_{1,1}v_{2,2} & - v_{1,3}v_{2,2} + v_{1,2}v_{2,3} \\
 - 2v_{1,2}v_{2,1} + 2v_{1,1}v_{2,2} & 0 & - v_{1,3}v_{2,1} + v_{1,1}v_{2,3} \\
v_{1,3}v_{2,2} - v_{1,2}v_{2,3} & v_{1,3}v_{2,1} - v_{1,1}v_{2,3} & 0
\end{pmatrix}\] Both blocks computed directly, not merely asserted from the recursive formula.

\end{example}

\begin{remark}{a false shortcut, caught and corrected}{auto9} A naive reading of \cref{:lem-jacrecursion} suggests appending \({\overline{v}}_{k} = 0\) preserves rank on an inherited component -- this is false: the old-variable block is left-multiplied by \(- \left\lbrack {\overline{v}}_{k} \right\rbrack_{\times}T\), not merely padded with zero rows. At the first inherited case (\(k = 2 \rightarrow 3\)), appending \({\overline{v}}_{3} = 0\) actually drops the rank to 0. The correct witness (\cref{:lem-witness-inherited} below) instead picks \({\overline{v}}_{k} = e_{i}\) for a standard basis vector \textbf{not} lying in \(T \cdot \text{ im }J_{k - 1}(q)\) -- a genuine choice, not a formality.

\end{remark}

\begin{lemma}{rank-2 witness on the new component}{auto10} For every \(k \geq 2\), there is a point \(x = (p,0)\), \(p\) the \cref{:lem-nondeg} non-degeneracy witness at level \(k - 1\), with \(x \in V\left( I_{k} \right)\) and \(\text{rank }J_{k(x)} = 2\).

\end{lemma} \protect\phantomsection\label{lem-witness-new}{}

\begin{proof}[Proof of \cref{:lem-witness-new}] By \cref{:lem-jacrecursion}, \(J_{k(p,0)} = \left\lbrack 0|\left\lbrack \tau({\overline{w}}_{k - 1}(p)) \right\rbrack_{\times} \right\rbrack\) since \(- \lbrack 0\rbrack_{\times}TJ_{k - 1}(p) = 0\); and \(\tau({\overline{w}}_{k - 1}(p)) \neq 0\) since \({\overline{w}}_{k - 1}(p) \neq 0\) (\cref{:lem-nondeg}) and \(\tau\) is invertible. For \(u \neq 0\), \(\ker\lbrack u\rbrack_{\times} = \left\{ v:u \times v = 0 \right\} = \mathbb{F}u\) (a cross product vanishes exactly on linearly dependent pairs), so \(\left\lbrack \tau({\overline{w}}_{k - 1}(p)) \right\rbrack_{\times}\) has rank 2, giving \(\text{rank }J_{k(p,0)} = 2\).

\end{proof}

\begin{lemma}{rank-2 witness on an inherited component}{lem-witness-inherited} For every \(k \geq 3\) and every component \(Z\) of \(V\left( I_{k - 1} \right)\) with a rank-2 witness \(q \in Z\) (\cref{:lem-witness-new}, or this lemma applied at \(k - 1\)), there is \(x = \left( q,e_{i} \right) \in Z \times \mathbb{A}^{3} \subset V\left( I_{k} \right)\) with \(\text{rank }J_{k(x)} = 2\).

\end{lemma} \protect\phantomsection\label{lem-witness-inherited}{}

\begin{proof}[Proof of \cref{:lem-witness-inherited}] Let \(W ≔ T \cdot \text{ im }J_{k - 1}(q)\), a 2-dimensional subspace of \(\mathbb{F}^{3}\) since \(T\) is invertible and \(\text{rank }J_{k - 1}(q) = 2\). Since \(e_{1},e_{2},e_{3}\) span \(\mathbb{F}^{3}\) and \(W \neq \mathbb{F}^{3}\), some \(e_{i} \notin W\); fix such an \(i\). Since \(q \in V\left( I_{k - 1} \right)\), \({\overline{w}}_{k - 1}(q) = 0\), so \cref{:lem-jacrecursion} gives \(J_{k\left( q,e_{i} \right)} = \left\lbrack - \left\lbrack e_{i} \right\rbrack_{\times}TJ_{k - 1}(q)|0 \right\rbrack\), of rank \(\dim(\left\lbrack e_{i} \right\rbrack_{\times}W)\) since \(TJ_{k - 1}(q)\) has image exactly \(W\). As \(\ker\left\lbrack e_{i} \right\rbrack_{\times} = \mathbb{F}e_{i}\) and \(W \cap \mathbb{F}e_{i} = 0\) (else \(e_{i} \in W\)), \(\left\lbrack e_{i} \right\rbrack_{\times}\) is injective on \(W\), so \(\dim(\left\lbrack e_{i} \right\rbrack_{\times}W) = 2\).

\end{proof}

\begin{theorem}{{}\(I_{k}\) is radical}{thm-radical} For every \(k \geq 2\): \(I_{k}\) is radical.

\end{theorem} \protect\phantomsection\label{thm-radical}{}

\begin{proof}[Proof of \cref{:thm-radical}] By induction on \(k\) (mirroring \cref{:thm-componentcount}`s own induction), \cref{:lem-witness-new} and \cref{:lem-witness-inherited} exhibit a rank-2-Jacobian point on \textbf{every} one of the \(k - 1\) components of \cref{:thm-componentcount}. \cref{:lem-jacobian-regularity} then gives \(R_{0}\) at every one (and, by its ``unique component'' clause, confirms each witness lies only on its intended component -- an independent check on \cref{:thm-componentcount}`s count, not merely assumed); \cref{:lem-serre} gives \(S_{1}\), using \cref{:thm-cm}`s Cohen-Macaulayness. Serre's criterion gives radicality.

\end{proof}

Verified independently in Sage at every \(k = 2,\ldots,6\): the Jacobian recursion (\cref{:lem-jacrecursion}) checked as a polynomial identity, not just at points, and each witness confirmed to lie on \textbf{exactly} its intended component via an independently-computed primary decomposition (code/verify\_jacobian\_rank\_all\_components.sage).

\section{The Primary Decomposition Theorem}

Every ingredient is now in hand for the fully general statement -- component count (§5) and radicality (§7) together, neither sufficient alone:

\begin{theorem}{Primary Decomposition Theorem}{thm-primarydecomp} For every \(k \geq 2\), over any field of characteristic \(\neq 2\): \(I_{k}\) has exactly \(k - 1\) irreducible, radical, primary components, namely \(P_{2},\ldots,P_{k}\) of \cref{:def-primecomponents}, and \[I_{k} = P_{2} \cap \cdots \cap P_{k}\] is an irredundant primary decomposition into \(k - 1\) primes.

\end{theorem} \protect\phantomsection\label{thm-primarydecomp}{}

\begin{proof}[Proof of \cref{:thm-primarydecomp}] Spelled out in full, combining the two independent halves proven above rather than left as a one-line conjunction. By \cref{:thm-radical}, \(I_{k}\) is radical, i.e. \(I_{k} = \sqrt{I_{k}}\). A standard fact about radical ideals in a Noetherian ring: a radical ideal equals the intersection of its own minimal primes, \(I_{k} = \cap_{P \in \text{ Min}\left( I_{k} \right)}P\) -- since \(\sqrt{I_{k}}\) is, by definition, the intersection of \textbf{all} primes containing \(I_{k}\), and the minimal ones already suffice (every non-minimal prime containing \(I_{k}\) contains some minimal one, hence is redundant in the intersection). \cref{:thm-componentcount} identifies \(\text{Min}\left( I_{k} \right)\) explicitly: exactly \(k - 1\) primes, one per block pattern \(j = 2,\ldots,k\), named \(P_{2},\ldots,P_{k}\) by \cref{:def-primecomponents}. Substituting gives the displayed equation. Each \(P_{j}\), being prime, is trivially its own primary-decomposition component (\(P_{j}\)-primary with associated prime \(P_{j}\) itself) -- no separate argument for ``primary'' beyond ``prime'' is needed, and primality of each \(P_{j}\) is exactly \cref{:lem-y2-irred} (applied at the appropriate parameter, via \cref{:thm-radical}`s own induction). Irredundancy -- no \(P_{j}\) can be omitted, i.e. no \(P_{j}\) contains \(\cap_{i \neq j}P_{i}\) -- is exactly \cref{:thm-componentcount}`s own no-containment conclusion (Proposition 17 and Proposition 18, applied at every level of that theorem's induction), not a fact needing any further proof here.

\end{proof}

\begin{remark}{summary}{auto11} \cref{:thm-codim} (codimension 2), \cref{:thm-mingens} (3 minimal generators), \cref{:thm-cm} (Cohen-Macaulay), and \cref{:thm-radical} (radical) together say exactly what a single bundled ``\(I_{k}\) has codimension 2, is radical, has exactly 3 minimal generators, and is Cohen-Macaulay'' statement would -- now as four independently-proven results, each established at the point in this article where its actual dependencies (and no others) were available, rather than one claim proven in parts across several sections.

\end{remark}

\section{Summary}

For every \(k \geq 2\), over any field \(\mathbb{F}\) of characteristic \(\neq 2\):

\begin{itemize}
\item
  \(I_{k}\) has codimension 2 (\cref{:thm-codim}), presented as the ideal of \(2 \times 2\) minors of an explicit \(2 \times 3\) matrix (\cref{:thm-determinantal}) -- a genuine determinantal ideal, not merely determinantal-looking.
\item
  \(I_{k}\) has exactly 3 minimal generators (\cref{:thm-mingens}), proven directly from linear independence, independent of the determinantal presentation.
\item
  \(R_{k}/I_{k}\) is Cohen-Macaulay (\cref{:thm-cm}), from the determinantal presentation and classical Eagon-Northcott/Hilbert-Burch theory (\cref{:prop-detideal}).
\item
  \(I_{k}\) is radical (\cref{:thm-radical}), via Serre's criterion: Cohen-Macaulayness gives \(S_{1}\) (\cref{:lem-serre}), and an explicit rank-2 Jacobian witness on every component gives \(R_{0}\) (\cref{:lem-witness-new}, \cref{:lem-witness-inherited}).
\item
  \(V\left( I_{k} \right)\) has exactly \(k - 1\) irreducible components (\cref{:thm-componentcount}), each rational -- in fact birational to an explicit affine parameter space (\cref{:prop-rational}) -- and these assemble, via the two facts above, into the irredundant primary decomposition of the Primary Decomposition Theorem (\cref{:thm-primarydecomp}): \(I_{k} = P_{2} \cap \cdots \cap P_{k}\).
\end{itemize}

Nothing above is conjectural or checked only at small \(k\): every statement is proven for \textbf{every} \(k \geq 2\), over an \textbf{arbitrary} field of the stated characteristic. The one running computational thread -- Sage verification of the determinantal identity, the Jacobian recursion, and the witness points, at \(k = 2,\ldots,6\) -- corroborates the general proofs; it is not a substitute for them anywhere in this article.

\section{References}

\begin{thebibliography}{99}
\bibitem{BrunsVetter1988} W. Bruns, U. Vetter, \emph{Determinantal Rings}, Lecture Notes in Mathematics 1327, Springer-Verlag, 1988. Out of print; freely available with Springer's permission at \href{https://www.home.uni-osnabrueck.de/wbruns/brunsw/detrings.pdf}{the first author's own page}. Theorem (2.7) and Theorem (16.36) verified directly against this copy, 2026-08-16.  \bibitem{EagonNorthcott1962} J.A. Eagon, D.G. Northcott, \emph{Ideals defined by matrices and a certain complex associated with them}, Proc. Roy. Soc. London Ser. A 269(1337) (1962), 188-204. DOI: \href{https://doi.org/10.1098/rspa.1962.0170}{10.1098/rspa.1962.0170}. Closed access (confirmed via Unpaywall, 2026-08-16); superseded here by Bruns-Vetter's more general treatment.  \bibitem{MotzkinTaussky1955} T. Motzkin, O. Taussky, \emph{Pairs of matrices with property L. II}, Trans. Amer. Math. Soc. 80 (1955), 387-401.  \bibitem{Stacks00OS} \emph{The Stacks Project}, tag 00OS (dimension formula for finitely generated field extensions / equicodimensionality of affine domains), and equivalently D. Eisenbud, \emph{Commutative Algebra with a View Toward Algebraic Geometry}, Springer GTM 150, 1995, Cor. 13.4. Open access at \href{https://stacks.math.columbia.edu/tag/00OS}{stacks.math.columbia.edu/tag/00OS}.  \bibitem{Gerstenhaber1961} M. Gerstenhaber, \emph{On dominance and varieties of commuting matrices}, Ann. of Math. 73 (1961), 324-348. JSTOR-gated; not freely available.  \bibitem{ArtinHochsterConjecture} M. Artin, M. Hochster, attributed conjecture: for \(I\) the ideal generated by the entries of \(XY - YX\), \(\text{Spec}(R/I)\) is reduced and Cohen-Macaulay. No single dateable original publication found; stated as ``Conjecture 1.1 (M. Artin, M. Hochster)'' in Majidi-Zolbanin and Snapp below.  \bibitem{MajidiZolbaninSnapp} M. Majidi-Zolbanin, B. Snapp, \emph{A note on the variety of pairs of matrices whose product is symmetric}, undated preprint. Open access at \href{https://people.math.osu.edu/snapp.14/SymmetrizingVariety.pdf}{people.math.osu.edu/ snapp.14/SymmetrizingVariety.pdf}.  \bibitem{Chau2024} T. Chau, \emph{The \(F\)-regularity of algebraic sets related to commutator matrices}, J. Algebra 652 (2024), 52-66. DOI: \href{https://doi.org/10.1016/j.jalgebra.2024.04.008}{10.1016/j.jalgebra.2024.04.008}. Closed access (confirmed via Unpaywall); obtained via Jan's LiU institutional access.  \bibitem{Hreinsdottir2005} F. Hreinsdóttir, \emph{Conjectures on the ring of commuting matrices}, arXiv:math/0501465. Open access at \href{https://arxiv.org/abs/math/0501465}{arxiv.org/abs/math/0501465}.  \bibitem{Knutson2003} A. Knutson, \emph{Some Schemes Related to the Commuting Variety}, arXiv:math/0306275 (to appear in J. Algebraic Geom.). Open access at \href{https://arxiv.org/abs/math/0306275}{arxiv.org/abs/math/0306275}.  \bibitem{Youcis2014} A. Youcis, answer to ``What can we say about the image of a regular map?'', Mathematics Stack Exchange, question 731178, answer 732415, 2014. States both the point-set-topology fact that the continuous image of an irreducible space is irreducible, and the \((x,y) \mapsto (x,xy)\) non-closed-image example cited above. CC BY-SA 3.0. Open access at \href{https://math.stackexchange.com/questions/731178/what-can-we-say-about-the-image-of-a-regular-map}{math.stackexchange.com/questions/731178}. 
\end{thebibliography}


Acknowledgements: the computational verification, literature search, and drafting of this article were carried out with substantial assistance from Claude (Anthropic), acting as an AI research assistant under the author's direction and review -- not as a co-author. Every claim above has been checked by the method stated -- symbolic/computational verification against the actual generating ideal, or citation checked directly against the cited primary source -- rather than presented on the strength of a pattern or a memory of the literature.

\end{document}
